\subsection{Mapa de Procesos}

El mapa de macroprocesos de la Universidad Nacional de Ingeniería organiza sus actividades en tres niveles interrelacionados para cumplir con su misión institucional y responder a las necesidades de la sociedad:

\begin{enumerate}
    \item \textbf{Procesos Estratégicos:} Conducen el rumbo de la universidad mediante el planeamiento estratégico, la asignación presupuestal, el aseguramiento continuo de la calidad y los convenios de cooperación interinstitucional.
    \item \textbf{Procesos Misionales (Operativos):} Constituyen la razón de ser de la institución. Incluyen los procesos de admisión, la gestión de la formación académica (pregrado y posgrado), la investigación científica y tecnológica, la proyección social y la graduación.
    \item \textbf{Procesos de Soporte (Apoyo):} Proveen los recursos y plataformas indispensables para la operación institucional, tales como la administración de registros académicos e infraestructura de tecnologías de información (DIRCE), los servicios de salud y bienestar estudiantil, la gestión financiera y administrativa (DIGA), y el mantenimiento físico del campus.
\end{enumerate}

\paragraph{Ubicación del proceso en estudio}
El \textbf{Procedimiento Administrativo de Matrícula Regular} se sitúa dentro del macroproceso misional de \textbf{Gestión de la Formación Académica}. Este proceso es el mecanismo formal mediante el cual el estudiante formaliza su compromiso académico para el ciclo lectivo, y requiere la coordinación directa con los procesos de soporte de registros académicos (\texttt{INTRALU}) y bienestar universitario (Centro Médico).

Para efectos del presente análisis y rediseño de procesos, la Matrícula Regular se desglosa en cuatro subprocesos secuenciales:
\begin{itemize}
    \item \textbf{Subproceso 1: Acreditación y Pre-requisitos.} Validación de cuentas de acceso, verificación de aportes obligatorios (Fondo de Autoseguro Estudiantil), cumplimiento de requisitos médicos y consulta del turno de matrícula.
    \item \textbf{Subproceso 2: Selección y Registro de Cursos.} Ingreso a la plataforma web institucional, consulta del catálogo de asignaturas y horarios disponibles de la facultad, y selección de asignaturas.
    \item \textbf{Subproceso 3.A: Verificación y Control de Reglas Académicas.} Validación informática en tiempo real de los requisitos curriculares (prerrequisitos de cursos, límite máximo de créditos y cruces de horario).
    \item \textbf{Subproceso 3.B: Confirmación y Registro Transaccional Oficial.} Grabación definitiva de la inscripción en la base de datos institucional, asignación formal de vacantes y emisión de la ficha de matrícula oficial.
\end{itemize}

En la Figura~\ref{fig:mapa_procesos_uni} se muestra el mapa de procesos de la universidad y el detalle del proceso de matrícula regular con sus cuatro subprocesos.

% ==============================================================================
% Diagrama: Mapa de Procesos Institucional de la UNI (Versión Compacta y Estable)
% ==============================================================================
\tikzset{
    cajaBase/.style={
        rectangle, rounded corners=3pt, thick, minimum height=1.15cm,
        align=center, font=\scriptsize\bfseries\sffamily
    },
    % #1: color | #2: ancho | #3: opciones extra
    caja/.style n args={3}{cajaBase, draw=#1!80!black, fill=#1!8, text width=#2, #3},
    % #1: color | #2: grosor
    flecha/.style 2 args={->, line width=#2, draw=#1!75!black},
    % Estilos puntuales por ID (activados automáticamente con /.try)
    mis2/.style={fill=teal!22, draw=teal!95!black},
    bandaSub/.style={densely dashed, line width=1.1pt}
}

% #1: args de flecha (opc) | #2: args de caja | #3: sep | #4: coord | #5: lista
\providecommand{\filaCajas}[5][]{%
    \foreach \id/\texto [count=\i, remember=\id as \prev] in {#5} {
        \tikzset{nodoActual/.style={caja=#2, \id/.try}}
        \ifnum\i=1
            \node[nodoActual] (\id) at (#4) {\texto};
        \else
            \node[nodoActual, right=#3 of \prev] (\id) {\texto};
            \if\relax\detokenize{#1}\relax\else
                \draw[flecha=#1] (\prev) -- (\id);
            \fi
        \fi
    }%
}

\begin{figure}[H]
\centering
\resizebox{\linewidth}{!}{%
\begin{tikzpicture}[>=Latex, font=\sffamily]

    % 1. FILAS DE PROCESOS (Estratégicos, Misionales, Subprocesos y Soporte)
    \filaCajas{{blue}{4.3cm}{font=\small\bfseries\sffamily}}{0.75cm}{-0.2, 7.6}{
        est1/{Planeamiento Institucional\\y Presupuesto},
        est2/{Gestión de la Calidad,\\Acreditación y Mejora},
        est3/{Dirección Universitaria\\y Cooperación Técnica}
    }

    \filaCajas[{teal}{1.4pt}]{{teal}{2.7cm}{}}{0.50cm}{-1.8, 4.8}{
        mis1/{Admisión\\e Ingreso},
        mis2/{Formación Integral\\(Pregrado y Posgrado)},
        mis3/{Investigación,\\Desarrollo e Innovación},
        mis4/{Responsabilidad Social\\y Extensión Cultural},
        mis5/{Graduación\\y Titulación}
    }

    \filaCajas[{red}{1.5pt}]{{red}{3.2cm}{fill=white, very thick, font=\scriptsize\sffamily}}{0.30cm}{-0.4, 2.15}{
        sub1/{\textbf{Subproceso 1}\\Acreditación y\\Pre-requisitos},
        sub2/{\textbf{Subproceso 2}\\Selección y\\Registro de Cursos},
        sub3a/{\textbf{Subproceso 3.A}\\Verificación de\\Reglas Académicas},
        sub3b/{\textbf{Subproceso 3.B}\\Confirmación y\\Registro Transaccional}
    }

    \draw[flecha={red}{1.2pt}, dashed] (mis2.south) -- (mis2.south |- 0, 3.35);

    \filaCajas{{purple}{3.3cm}{}}{0.35cm}{-0.7, -0.5}{
        sop1/{Registros Académicos\\y TI (DIRCE / Intralu)},
        sop2/{Bienestar y Salud\\(Centro Médico / Autoseguro)},
        sop3/{Gestión Administrativa\\y Financiera (DIGA)},
        sop4/{Mantenimiento,\\Infraestructura y Logística}
    }

    % 2. CAPAS DE FONDO Y CONTENEDORES (Unificadas con \id/.try)
    \begin{scope}[on background layer]
        \foreach \id/\col/\xA/\yA/\xB/\yB/\tit in {
            bandaEst/blue/-3.7/6.7/14.0/8.9/PROCESOS ESTRATÉGICOS,
            bandaMis/teal/-3.7/1.25/14.0/6.1/PROCESOS MISIONALES (OPERATIVOS),
            bandaSub/red/-2.3/1.35/12.8/3.35/PROCESO EN ESTUDIO: MATRÍCULA REGULAR (Facultad de Ciencias),
            bandaSop/purple/-3.7/-1.4/14.0/0.8/PROCESOS DE SOPORTE%
        }{
            \filldraw[rounded corners=5pt, fill=\col!3, draw=\col!45!black, thick, \id/.try] (\xA,\yA) rectangle (\xB,\yB);
            \node[fill=\col!75!black, text=white, font=\bfseries\footnotesize\sffamily, rounded corners=2pt, 
                  anchor=north west, inner xsep=6pt, inner ysep=3pt] at (\xA,\yB) {\tit};
        }
    \end{scope}

    % 3. FLANCOS LATERALES Y FLECHAS DE FLUJO
    \foreach \id/\x/\ang/\xF/\xT/\tit/\sub in {
        entrada/-5.1/90/-4.1/-3.7/NECESIDADES Y EXPECTATIVAS/{Requisitos de Estudiantes, Comunidad y Sector Productivo},
        salida/15.4/-90/14.0/14.4/SATISFACCIÓN Y RESULTADOS/{Profesionales Graduados, Innovación y Aporte al País}%
    }{
        \node[cajaBase, rounded corners=4pt, draw=black!70, fill=gray!18, line width=1.1pt, minimum width=1.8cm, minimum height=10.6cm] 
            (\id) at (\x, 3.75) {\rotatebox{\ang}{\parbox{9.5cm}{\centering\bfseries\small\tit\\[0.1cm]\footnotesize\normalfont\sub}}};
        \draw[flecha={teal}{3pt}] (\xF, 4.8) -- (\xT, 4.8);
    }

\end{tikzpicture}%
}
\caption{Mapa de procesos institucional de la Universidad Nacional de Ingeniería y ubicación del proceso de matrícula regular.}
\label{fig:mapa_procesos_uni}
\end{figure}
